\subsubsection{Interaktives Benuterzeroberflächen Design und Abfragelogik}

Die GUI des Prototyps ist funktional in zwei Hauptbereiche unterteilt. Der zentrale Zone bildet der 3D-Viewport, in welchem über Maussteuerung navigiert werden kann. Auf der rechten Seite befindet sich eine Leiste, die als Informations- und Steuerungspaneel fungiert \cite{wangAutomatedPhenotypicTrait2024}. Der obere Teil ist der Anzeige gesammelter Baumkennwerte nach einer Selektion vorbehalten. Der untere Abschnitt enthält die interaktiven Visualisierungseinstellungen, über die Nutzer die Anzeige des DOP, der Kronenpolygone, der Farbskala sowie des Tiefenschattens wahlweise aktivieren oder deaktivieren können. Ein integriertes Dropdown-Menü ermöglicht zudem die Farbkodierung der Punktwolke auf Basis unterschiedlicher Attribute. Die Standardeinstellung färbt die Segmente nach ihrer eindeutigen ID ein, um benachbarte Individuen visuell voneinander abzugrenzen. Alternativ stehen Einfärbungen nach dem NDVI, der absoluten Höhe sowie der Kronenhöhe zur Verfügung. Darüber hinaus lässt sich der Status des Katasterabgleichs visualisieren. Bei einem Match werden die Katasterattribute automatisch im Informationsbereich mit angezeigt \cite{wegenNonPhotorealisticRendering3D2024}.

Die informationstechnische Verknüpfung zwischen der 3D-Punktwolke und den zugeordneten Informationen wird durch eine interaktive Selektionslogik mittels 3D-to-World-Raycasting realisiert. Bei der Auswahl eines Objekts im Viewport wird die 2D-Bildschirmkoordinate der Maus als virtueller Strahl in den Objektraum projiziert, um den Schnittpunkt mit der Basisebene des Geländes zu berechnen. An dieser Position erfolgt eine räumliche Distanzabfrage gegen die Geometrie der Kronenpolygone \cite{kempfObliqueViewIndividual2021}. Bei einem positiven Treffer werden die mit der Baum-ID verknüpften LiDAR-Daten sowie die Katasterattribute ausgelesen und in der Seitenleiste angezeigt. 